% 无限测度的反例：I_k=(2^{-k},2^{-k+1}]，lambda(I_k)=2^{-k}，nu=sum_k 2^k lambda|_{I_k}给出nu(I_k)=1；画出k=1,...,4，横轴1单位=9cm。
\begin{tikzpicture}[
  lab/.style={font=\figurefont\fontsize{8}{10}\selectfont,text=FigureInk},
  note/.style={font=\FigureNoteFont\fontsize{7.5}{9}\selectfont,text=FigureMuted,align=center},
  ax/.style={draw=FigureStroke,line width=.75pt}]
\def\u{9}
\fill[FigureBlueFill] ({\u/2},0) rectangle ({\u},.35);
\fill[FigureGreenFill] ({\u/4},0) rectangle ({\u/2},.35);
\fill[FigureBlueFill] ({\u/8},0) rectangle ({\u/4},.35);
\fill[FigureGreenFill] ({\u/16},0) rectangle ({\u/8},.35);
\fill[FigureGrid] (0,0) rectangle ({\u/16},.35);
\draw[ax] (0,0)--(\u,0);
\foreach \t/\s in {1/1,0.5/\frac12,0.25/\frac14,0.125/\frac18}{
  \draw[FigureStroke,line width=.6pt] ({\u*\t},0)--({\u*\t},.45);
  \node[lab,below] at ({\u*\t},-.02) {$\s$};}
\draw[FigureStroke,line width=.6pt] (0,0)--(0,.45); \node[lab,below] at (0,-.02) {$0$};
% labels above
\foreach \k/\a/\b/\l in {1/0.5/1/2^{-1},2/0.25/0.5/2^{-2},3/0.125/0.25/2^{-3}}{
  \node[lab] at ({\u*(\a+\b)/2},.85) {$I_{\k}$};
  \node[note] at ({\u*(\a+\b)/2},1.35) {$\lambda(I_{\k})=\l$};
  \node[note] at ({\u*(\a+\b)/2},1.75) {$\nu(I_{\k})=1$};}
\node[lab] at ({\u*3/32},.85) {$I_4$};
\node[note,anchor=west] at ({\u*3/32-.3},2.25) {$\lambda(I_k)\to0$，$\nu(I_k)=1$};
\draw[FigureMuted,line width=.5pt] ({\u*3/32},1.05)--({\u*3/32},2.05);
\end{tikzpicture}
